\begin{ouverture}[espace]{18}{Solides, patrons et espace}{Partie 4 · Espace et géométrie}
{Cubes · Pavés · Patrons · Assemblages · Déplacements}
\objectif{Nommer et décrire les solides usuels.}
\objectif{Reconnaître et construire un patron du cube.}
\objectif{Dénombrer des cubes et coder des déplacements.}
\prerequis{Figures planes usuelles et quadrillage.}
\end{ouverture}
\pagestyle{otMain}\essentiel\fontsize{11}{15}\selectfont
\phantomsection\label{cours-18-1}
\essSection{1}{Reconnaître les solides}
\begin{definition}Les solides usuels sont le cube, le pavé (ou parallélépipède rectangle), la boule, le cylindre, le cône, la pyramide et le prisme droit. Les solides dont toutes les faces sont plates sont des polyèdres ; la boule, le cylindre et le cône ont des surfaces courbes.\end{definition}
\label{fig-cours-18-1-cours-solides}\figureGeo{cours-solides}

\begin{propriete}Une face est une surface plate du solide ; une arête est un segment où deux faces se rejoignent ; un sommet est un point où des arêtes se rencontrent.\end{propriete}
\begin{attention}En perspective, certaines faces, arêtes et sommets sont cachés : les arêtes invisibles se dessinent en pointillés. Les faces carrées peuvent sembler penchées : la perspective déforme les angles, pas la réalité.\end{attention}
\clearpage
\phantomsection\label{cours-18-2}
\essSection{2}{Décrire les solides}
\begin{propriete}Le cube a 6 faces carrées, 12 arêtes et 8 sommets. Le pavé a 6 faces rectangulaires, 12 arêtes et 8 sommets : ses faces opposées sont identiques.\end{propriete}
\begin{propriete}Une pyramide a une base polygonale et des faces latérales triangulaires qui se rejoignent en un sommet. Un prisme droit a deux bases superposables et parallèles, reliées par des faces rectangulaires.\end{propriete}
\begin{exemple}Décrire, c’est donner le nombre et la nature des faces : « 6 faces carrées » suffit à reconnaître le cube ; « deux bases triangulaires et trois faces rectangulaires » décrit un prisme droit à base triangulaire.\end{exemple}
\begin{savoirfaire}[Identifier en perspective]Pour identifier un solide en perspective : compter les faces visibles et cachées, observer la forme des faces, chercher les bases. Justifier par les propriétés, pas par l’allure.\end{savoirfaire}
\clearpage
\phantomsection\label{cours-18-3}
\essSection{3}{Le patron du cube}
\begin{definition}Un patron d’un solide est un dessin d’un seul tenant qui, plié le long des arêtes, reconstruit le solide. Un patron du cube est formé de 6 carrés identiques bien disposés.\end{definition}
\label{fig-cours-18-3-cours-patron}\figureGeo{cours-patron}

\begin{savoirfaire}[Plier mentalement]Pour tester un patron de cube, plier mentalement : choisir une face de fond, puis relever les autres une à une. Si deux carrés se superposent ou si une face manque, ce n’est pas un patron.\end{savoirfaire}
\begin{exemple}La croix de 6 carrés est un patron du cube. Quatre carrés alignés avec deux carrés du même côté n’en forment pas un : deux faces se superposent. Sur un patron du cube, les faces opposées sont séparées par une face.\end{exemple}
\clearpage
\phantomsection\label{cours-18-4}
\essSection{4}{Assemblages et déplacements}
\begin{savoirfaire}[Compter les cubes cachés]Pour compter les cubes d’un assemblage, compter étage par étage sans oublier les cubes cachés qui portent les étages supérieurs. Un cube en hauteur repose toujours sur un cube ou un empilement.\end{savoirfaire}
\begin{exemple}Un déplacement sur quadrillage se code avec un vocabulaire précis : avancer d’un nombre de cases, pivoter d’un quart de tour à droite ou à gauche. La même suite d’instructions, exécutée depuis un autre départ, donne un autre chemin.\end{exemple}
\begin{savoirfaire}[Coder et exécuter]Pour exécuter un programme de déplacement : marquer la case et l’orientation de départ, puis suivre les instructions une à une en dessinant le chemin. Pour coder un chemin : le découper en avancées et en quarts de tour.\end{savoirfaire}
\begin{attention}« À droite » dépend de l’orientation de celui qui avance : après un demi-tour, sa droite a changé de côté. Toujours suivre l’orientation du personnage, pas celle du lecteur.\end{attention}

\clearpage\rubrique{Découverte}
\begin{blocLivre}{ex-18-1}\begin{exercice}[1]{Nommer les solides}
Nommer chacun des six solides dessinés.
\label{fig-ex-18-1}\figureGeo{sol-1}
\end{exercice}\end{blocLivre}
\begin{blocLivre}{ex-18-2}\begin{exercice}[1]{La carte d’identité du cube}
Donner le nombre de faces, d’arêtes et de sommets d’un cube, et la nature de ses faces.
\label{fig-ex-18-2}\figureGeo{sol-2}
\end{exercice}\end{blocLivre}
\begin{blocLivre}{ex-18-3}\begin{exercice}[1]{Décrire un pavé}
Décrire le pavé dessiné : nombre et nature des faces, arêtes et sommets. Que dire des faces opposées ?
\label{fig-ex-18-3}\figureGeo{sol-3}
\end{exercice}\end{blocLivre}
\begin{blocLivre}{ex-18-4}\begin{exercice}[1]{Les faces de la pyramide}
La pyramide dessinée a une base carrée. Donner la nature et le nombre de ses faces.
\label{fig-ex-18-4}\figureGeo{sol-4}
\end{exercice}\end{blocLivre}
\begin{blocLivre}{ex-18-5}\begin{exercice}[1]{Le prisme droit}
Décrire le prisme droit dessiné : ses bases et ses autres faces.
\label{fig-ex-18-5}\figureGeo{sol-5}
\end{exercice}\end{blocLivre}

\rubrique{Entraînement}
\begin{blocLivre}{ex-18-6}\begin{exercice}[2]{Patron ou pas ?}
Parmi les trois dispositions de six carrés, lesquelles sont des patrons de cube ?
\label{fig-ex-18-6}\figureGeo{sol-6}
\end{exercice}\end{blocLivre}
\begin{blocLivre}{ex-18-7}\begin{exercice}[2]{Les faces opposées}
Sur le patron en croix, la face du milieu est posée au fond. Quelle face se retrouve opposée à la face marquée A ?
\label{fig-ex-18-7}\figureGeo{sol-7}
\end{exercice}\end{blocLivre}
\begin{blocLivre}{ex-18-8}\begin{exercice}[2]{Compléter le patron}
Cinq carrés sont dessinés ; il manque une face. Donner une position du sixième carré qui donne un patron de cube, et une position qui n’en donne pas. Justifier.
\label{fig-ex-18-8}\figureGeo{sol-8}
\end{exercice}\end{blocLivre}
\begin{blocLivre}{ex-18-9}\begin{exercice}[2]{Les cubes cachés}
Compter les cubes de l’assemblage dessiné, en n’oubliant pas les cubes cachés.
\label{fig-ex-18-9}\figureGeo{sol-9}
\end{exercice}\end{blocLivre}
\begin{blocLivre}{ex-18-10}\begin{exercice}[2]{La bonne vue}
Associer à l’assemblage la vue de face et la vue de dessus parmi les trois vignettes proposées.
\label{fig-ex-18-10}\figureGeo{sol-10}
\end{exercice}\end{blocLivre}

\rubrique{Problèmes}
\begin{blocLivre}{ex-18-11}\begin{exercice}[3]{Construire un patron}
Tracer un patron d’un cube de 3 cm d’arête, en vraie grandeur.
\label{fig-ex-18-11}\figureGeo{sol-11}
\end{exercice}\end{blocLivre}
\begin{blocLivre}{ex-18-12}\begin{exercice}[3]{Agrandir l’assemblage}
L’assemblage dessiné compte 7 cubes. Combien de cubes faut-il ajouter pour obtenir un grand cube de 2 cubes de côté ? Et pour un grand cube de 3 cubes de côté ?
\label{fig-ex-18-12}\figureGeo{sol-12}
\end{exercice}\end{blocLivre}
\begin{blocLivre}{ex-18-13}\begin{exercice}[3]{Coder un chemin}
Écrire la suite d’instructions qui mène le pion de la case départ à la case but, en suivant le chemin dessiné : avancer, pivoter à droite ou à gauche.
\label{fig-ex-18-13}\figureGeo{sol-13}
\end{exercice}\end{blocLivre}
\begin{blocLivre}{ex-18-14}\begin{exercice}[3]{Exécuter un programme}
Le pion part de la case marquée, tourné vers la droite de la feuille. Exécuter : avancer de 2 cases ; pivoter d’un quart de tour à gauche ; avancer de 2 cases ; pivoter d’un quart de tour à gauche ; avancer de 1 case. Où arrive-t-il ?
\label{fig-ex-18-14}\figureGeo{sol-14}
\end{exercice}\end{blocLivre}
\begin{blocLivre}{ex-18-15}\begin{exercice}[3]{Le dé à compléter}
Sur un dé, les points des faces opposées totalisent toujours 7. Compléter le patron du dé : les faces 1, 2 et 3 sont déjà placées. Où écrire 4, 5 et 6 ?
\label{fig-ex-18-15}\figureGeo{sol-15}
\end{exercice}\end{blocLivre}

\clearpage
